\chapter{Data Model and Persistence}\label{ch:ch05}

Every exam that a teacher authors, every class a student joins, every line
of code a student submits, and every verdict that Judge0 returns is, in the
end, a row in a PostgreSQL~16 table~\cite{postgres}. This chapter is the map
of that store. It presents the fifteen entities that make up the APEX schema,
walks the columns and constraints that matter, explains the three composite
unique indexes that carry the system's correctness invariants, and then
turns to how the schema changes over time --- the migration pipeline and the
one incident that taught the team why that pipeline exists at all.

The persistence layer is deliberately small. APEX is a single Go binary
talking to a single database through a single object-relational mapper
(GORM~\cite{gorm}); there is no sharding, no read replica, no separate
analytics store. The design goal was a schema that a single engineer can
hold in their head, on the reasoning that a schema you cannot trace by hand
is a schema you will eventually corrupt. The rest of this chapter argues that
this smallness is a feature: it is what makes the grading pipeline's
correctness auditable, and it is what let the team recover cleanly from a
migration mistake that would have been catastrophic in a larger design.

\section{The Entity Landscape}

The canonical roster of entities is the argument list of GORM's
\texttt{AutoMigrate} call in \texttt{internal/database/database.go}. Fifteen
Go structs are registered there, and each maps to exactly one table:

\begin{lstlisting}[language=Go,caption={The fifteen registered models, in migration order (internal/database/database.go).},label={lst:automigrate}]
DB.AutoMigrate(
    &models.User{},
    &models.UserProfile{},
    &models.Class{},
    &models.ClassMember{},
    &models.Exam{},
    &models.ExamClass{},
    &models.Problem{},
    &models.TestCase{},
    &models.ExamAttempt{},
    &models.Submission{},
    &models.TestResult{},
    &models.Notification{},
    &models.Announcement{},
    &models.Folder{},
    &models.PasswordResetToken{},
)
\end{lstlisting}

These fifteen tables cluster into four \emph{aggregates} --- a
domain-driven-design term for a group of entities treated as one consistency
unit with a single root that owns the rest~\cite{fowler2002patterns,
kleppmann2017ddia}. \Cref{tab:aggregates} lays out the grouping. The counts
form a convenient mnemonic: identity has three tables, classes four,
assessment seven, and one cross-cutting table stands alone, for a total of
$3+4+7+1=15$.

\begin{table}[htbp]
\centering
\caption{The four aggregates of the APEX data model.}
\label{tab:aggregates}
\begin{tabularx}{\linewidth}{@{}L{2.6cm} L{3.0cm} X@{}}
\toprule
\textbf{Aggregate} & \textbf{Root} & \textbf{Members} \\
\midrule
Identity (3)   & \texttt{User}        & \texttt{UserProfile} (1:1), \texttt{PasswordResetToken} (1:N) \\
Classes (4)    & \texttt{Class}       & \texttt{ClassMember} (M:N with \texttt{User}), \texttt{Announcement} (1:N), \texttt{Folder} (teacher-owned) \\
Assessment (7) & \texttt{Exam}        & \texttt{ExamClass} (M:N with \texttt{Class}), \texttt{Problem} (1:N) $\rightarrow$ \texttt{TestCase} (1:N); \texttt{ExamAttempt} (1:N) $\rightarrow$ \texttt{Submission} (1:N) $\rightarrow$ \texttt{TestResult} (1:N) \\
System (1)     & ---                  & \texttt{Notification} (belongs to a \texttt{User}, cross-cutting) \\
\bottomrule
\end{tabularx}
\end{table}

The assessment cluster is large enough that it is often useful to split it
further into an \emph{authoring} half (\texttt{Exam}, \texttt{ExamClass},
\texttt{Problem}, \texttt{TestCase} --- the entities a teacher builds) and a
\emph{taking} half (\texttt{ExamAttempt}, \texttt{Submission},
\texttt{TestResult} --- the entities a student produces). The taking half is
the one governed by the system's most important database constraint, and
\cref{sec:attempt-aggregate} gives it its own treatment.
\Cref{fig:er} shows the complete entity-relationship diagram: the four
aggregates, the foreign-key edges between them, and the join tables that
model the two many-to-many relationships.

\begin{figure}[htbp]\centering
\resizebox{\textwidth}{!}{%
\begin{tikzpicture}[
  every node/.style={font=\scriptsize},
  entity/.style={draw=apexblue!85, fill=apexlight, rounded corners=1.5pt,
    rectangle split, rectangle split parts=2, rectangle split part align=left,
    align=left, inner sep=2.6pt, text width=25mm,
    rectangle split part fill={apexblue!14,apexlight}},
  jtable/.style={draw=apexgreen!70!black, fill=apexgreen!8, rounded corners=1.5pt,
    rectangle split, rectangle split parts=2, rectangle split part align=left,
    align=left, inner sep=2.6pt, text width=25mm,
    rectangle split part fill={apexgreen!22,apexgreen!7}},
  uxnote/.style={draw=apexamber, fill=apexamber!14, rounded corners=1pt,
    font=\tiny, inner sep=2pt, align=left, text width=27mm},
  card/.style={font=\tiny\bfseries, text=apexblue!75!black, fill=white,
    inner sep=0.8pt},
  rel/.style={-{Latex[length=1.8mm]}, draw=apexgrey, semithick},
  jrel/.style={-{Latex[length=1.8mm]}, draw=apexgreen!60!black, semithick},
  opt/.style={-{Latex[length=1.8mm]}, draw=apexamber!85!black, semithick, dashed},
]

% ---- Entities: User cluster (left) ----
\node[entity] (usr)  at (0,2.6)
  {\textbf{User}\nodepart{two}\underline{id} PK\\ email (U)\\ role\\ password\_hash};
\node[entity] (prof) at (-3.6,5.0)
  {\textbf{UserProfile}\nodepart{two}\underline{id} PK\\ user\_id (U)\\ bio, avatar\_url\\ notify\_*};
\node[entity] (notf) at (-3.6,2.6)
  {\textbf{Notification}\nodepart{two}\underline{id} PK\\ user\_id (FK)\\ type, title\\ read};
\node[entity] (prst) at (-3.6,0.2)
  {\textbf{PasswordResetToken}\nodepart{two}\underline{id} PK\\ user\_id (FK)\\ token\_hash (U)\\ expires\_at};

% ---- Classes / enrolment ----
\node[entity] (cls)  at (3.9,5.2)
  {\textbf{Class}\nodepart{two}\underline{id} PK\\ teacher\_id (FK)\\ invite\_code (U)\\ passing\_threshold};
\node[jtable] (cmem) at (3.9,2.6)
  {\textbf{ClassMember}\nodepart{two}class\_id (FK)\\ user\_id (FK)\\ joined\_at};
\node[entity] (fold) at (3.9,-1.4)
  {\textbf{Folder}\nodepart{two}\underline{id} PK\\ teacher\_id (FK)\\ name};

% ---- Exams ----
\node[entity] (ann)  at (7.8,7.4)
  {\textbf{Announcement}\nodepart{two}\underline{id} PK\\ class\_id (FK)\\ teacher\_id (FK)\\ title, body};
\node[entity] (exam) at (7.8,5.2)
  {\textbf{Exam}\nodepart{two}\underline{id} PK\\ teacher\_id (FK)\\ title\\ is\_draft, start\_time};
\node[jtable] (exc)  at (7.8,2.9)
  {\textbf{ExamClass}\nodepart{two}exam\_id (FK)\\ class\_id (FK)};
\node[jtable] (att)  at (7.8,0.2)
  {\textbf{ExamAttempt}\nodepart{two}\underline{id} PK\\ user\_id (FK)\\ exam\_id (FK)\\ score, status};

% ---- Problems / submissions ----
\node[entity] (prob) at (11.7,5.2)
  {\textbf{Problem}\nodepart{two}\underline{id} PK\\ exam\_id (FK?)\\ folder\_id (FK?)\\ type, points};
\node[entity] (subm) at (11.7,0.2)
  {\textbf{Submission}\nodepart{two}\underline{id} PK\\ exam\_attempt\_id (FK?)\\ problem\_id (FK)\\ status, score};

% ---- Test cases / results ----
\node[entity] (tc)   at (15.6,5.2)
  {\textbf{TestCase}\nodepart{two}\underline{id} PK\\ problem\_id (FK)\\ input, expected\\ is\_sample};
\node[entity] (tr)   at (15.6,0.2)
  {\textbf{TestResult}\nodepart{two}\underline{id} PK\\ submission\_id (FK)\\ test\_case\_id (FK)\\ passed};

% ---- Relationships (cardinalities: 1 near parent, * near child) ----
\draw[rel] (usr) -- node[card,near start]{1} node[card,near end]{1} (prof);
\draw[rel] (usr) -- node[card,near start]{1} node[card,near end]{*} (notf);
\draw[rel] (usr) -- node[card,near start]{1} node[card,near end]{*} (prst);

\draw[rel] (usr) -- node[card,pos=0.12]{1} node[card,pos=0.9]{*} (cls);
\draw[jrel] (usr) -- node[card,pos=0.12]{1} node[card,pos=0.9]{*} (cmem);
\draw[jrel] (cls) -- node[card,near start]{1} node[card,near end]{*} (cmem);

\draw[rel] (usr.north east) to[out=25,in=200] node[card,pos=0.1]{1} node[card,pos=0.9]{*} (exam.west);
\draw[jrel] (exam) -- node[card,near start]{1} node[card,near end]{*} (exc);
\draw[jrel] (cls) -- node[card,pos=0.15]{1} node[card,pos=0.92]{*} (exc);

\draw[rel] (cls.north) to[out=70,in=180] node[card,pos=0.15]{1} node[card,pos=0.9]{*} (ann.west);

\draw[rel] (usr.south east) to[out=-30,in=205] node[card,pos=0.08]{1} node[card,pos=0.93]{*} (att.west);
\draw[jrel] (exam) -- node[card,near start]{1} node[card,near end]{*} (att);

\draw[rel] (usr.south) to[out=-70,in=180] node[card,pos=0.12]{1} node[card,pos=0.9]{*} (fold.west);
\draw[rel] (exam) -- node[card,near start]{1} node[card,near end]{*} (prob);
\draw[opt]  (fold.east) to[out=-8,in=250] node[card,pos=0.08]{1} node[card,pos=0.95]{*} (prob.south west);

\draw[rel] (prob) -- node[card,near start]{1} node[card,near end]{*} (tc);

\draw[rel] (att) -- node[card,near start]{1} node[card,near end]{*} (subm);
\draw[rel] (prob) -- node[card,near start]{1} node[card,near end]{*} (subm);

\draw[rel] (subm) -- node[card,near start]{1} node[card,near end]{*} (tr);
\draw[rel] (tc) -- node[card,near start]{1} node[card,near end]{*} (tr);

% ---- Unique composite-index notes ----
\node[uxnote,below=2.2mm of cmem] (n1)
  {\textbf{idx\_class\_user} (UNIQUE)\\ ClassMember(class\_id, user\_id)};
\node[uxnote,right=2.0mm of exc] (n2)
  {\textbf{idx\_exam\_class} (UNIQUE)\\ ExamClass(exam\_id, class\_id)};
\node[uxnote,below=2.2mm of att] (n3)
  {\textbf{idx\_user\_exam\_attempt} (UNIQUE)\\ ExamAttempt(user\_id, exam\_id)};
\draw[apexamber,densely dotted] (n1) -- (cmem);
\draw[apexamber,densely dotted] (n2) -- (exc);
\draw[apexamber,densely dotted] (n3) -- (att);

% ---- Legend ----
\node[draw=apexgrey!60, rounded corners=1.5pt, fill=white, align=left,
      font=\tiny, inner sep=3pt, text width=44mm] at (13.6,7.6)
  {\textbf{Legend}\\[1pt]
   \tikz{\draw[rel] (0,0)--(0.5,0);}~1..* relationship \quad
   \tikz{\draw[jrel] (0,0)--(0.5,0);}~join table\\
   \tikz{\draw[opt] (0,0)--(0.5,0);}~optional (nullable) FK \quad
   (U)~unique\quad (PK)~primary key};

\end{tikzpicture}%
}
\caption{Entity-relationship diagram of the APEX data model.}\label{fig:er}
\end{figure}


\section{Identity Tables}

The identity aggregate stores who may log in and how. Its root is
\texttt{User}:

\begin{lstlisting}[language=Go,caption={The User model (internal/models/user.go).},label={lst:user}]
type User struct {
    ID           uint      `gorm:"primaryKey" json:"id"`
    Email        string    `gorm:"size:255;uniqueIndex;not null" json:"email"`
    PasswordHash string    `gorm:"type:text;not null" json:"-"`
    Role         string    `gorm:"size:20;not null" json:"role"`
    Name         string    `gorm:"size:255;not null" json:"name"`
    GoogleID     string    `gorm:"size:64;index" json:"-"`
    CreatedAt    time.Time `json:"created_at"`
}
\end{lstlisting}

Several conventions visible here recur across every table. The primary key is
always a field named \texttt{ID} of Go type \texttt{uint} carrying
\texttt{gorm:"primaryKey"}, which PostgreSQL materializes as a
\texttt{BIGSERIAL} auto-incrementing column. String columns declare an
explicit \texttt{size:} so they become bounded \texttt{VARCHAR(n)} rather
than unbounded \texttt{text}; free-form fields such as \texttt{PasswordHash}
opt into \texttt{type:text} instead. The \texttt{json:"-"} tag on
\texttt{PasswordHash} and \texttt{GoogleID} instructs the marshaller never to
serialize those fields into an API response, so a password hash cannot leak
through a handler that carelessly returns a whole \texttt{User}. The
\texttt{Role} column holds one of exactly two values, \texttt{teacher} or
\texttt{student}; APEX has no administrator role.

\texttt{UserProfile} holds per-user preferences in a 1:1 relationship with
\texttt{User}, enforced by \texttt{UserID gorm:"uniqueIndex"} --- because the
foreign key is unique, a user has at most one profile row. It carries the
notification toggles (\texttt{notify\_email}, \texttt{notify\_push},
\texttt{notify\_exam\_reminders}, \texttt{notify\_results},
\texttt{notify\_exam\_email}) that the reminder scheduler consults before
fanning out, plus a small block of per-teacher creation defaults
(\texttt{default\_exam\_draft}, \texttt{default\_passing\_threshold},
\texttt{block\_announce\_with\_pending}). The distinction between these
defaults and column defaults is central to the incident discussed in
\cref{sec:incident}: \texttt{default\_exam\_draft} seeds a \emph{new} exam's
draft state in application code at creation time, and is not the same thing
as a DDL default backfilling existing rows.

\texttt{PasswordResetToken} completes the aggregate. It never stores the raw
reset token; only its SHA-256 hash lives in \texttt{token\_hash}
(\texttt{size:64;uniqueIndex;not null}, hidden from JSON), so a database leak
hands an attacker no usable reset links. The nullable \texttt{used\_at
*time.Time} makes the token single-use: a non-null value means ``already
spent,'' and together with \texttt{expires\_at} it forms a one-shot,
time-limited credential.

\section{Class Tables}

The class aggregate models enrolment and classroom communication. A
\texttt{Class} belongs to one teacher (\texttt{TeacherID}, with a
\texttt{Teacher User} belongs-to association loaded on demand via
\texttt{Preload}) and is joined through an eight-character
\texttt{invite\_code} (\texttt{size:8;uniqueIndex;not null}). The
\texttt{grades\_announced} boolean is the gate that releases results to
students: until a teacher flips it on, graded scores stay hidden.

Enrolment is a many-to-many relationship between \texttt{User} and
\texttt{Class}, and APEX models it with an explicit join \emph{entity},
\texttt{ClassMember}, rather than GORM's implicit \texttt{many2many} tag:

\begin{lstlisting}[language=Go,caption={The ClassMember join entity (internal/models/class.go).},label={lst:classmember}]
type ClassMember struct {
    ID       uint      `gorm:"primaryKey" json:"id"`
    ClassID  uint      `gorm:"not null;index;uniqueIndex:idx_class_user" json:"class_id"`
    Class    Class     `gorm:"foreignKey:ClassID" json:"-"`
    UserID   uint      `gorm:"not null;index;uniqueIndex:idx_class_user" json:"user_id"`
    User     User      `gorm:"foreignKey:UserID" json:"user,omitempty"`
    JoinedAt time.Time `gorm:"autoCreateTime" json:"joined_at"`
}
\end{lstlisting}

The explicit join entity is chosen so the relationship can carry its own
data --- here a \texttt{joined\_at} timestamp stamped automatically by
\texttt{autoCreateTime} --- and its own named constraint. That constraint is
the first of the three composite unique indexes: \texttt{idx\_class\_user}
appears on \emph{both} \texttt{ClassID} and \texttt{UserID} under the same
index name, which GORM compiles into a single unique index over the pair
\texttt{(class\_id, user\_id)}. Its effect is that a student physically
cannot join the same class twice --- the database rejects a duplicate pair,
independent of any application-level check.

The aggregate is completed by \texttt{Announcement}, which belongs to a class
and a teacher and stores an \texttt{attachments} list as a JSONB column (see
\cref{sec:jsonb}); posting one fans out \texttt{Notification} rows to every
enrolled student. \texttt{Folder} is a teacher-owned grouping for bank
questions, referenced by \texttt{Problem.folder\_id} with
\texttt{ON DELETE SET NULL} semantics so that deleting a folder unfiles its
questions rather than destroying them.

\section{Assessment Tables and the Polymorphic Problem}

The authoring half of the assessment aggregate begins with \texttt{Exam},
which owns many \texttt{Problem} rows and is assigned to classes through the
\texttt{ExamClass} join entity. \texttt{ExamClass} carries the second
composite unique index, \texttt{idx\_exam\_class} over \texttt{(exam\_id,
class\_id)}, forbidding the same exam from being assigned to the same class
twice. Its scheduling columns (\texttt{start\_time}, \texttt{end\_time},
\texttt{reset\_at}, and three reminder-dedup timestamps) are all
\texttt{*time.Time} pointers, so they map to nullable \texttt{TIMESTAMPTZ}
columns; a null \texttt{start\_time} means ``not scheduled yet,'' and an exam
cannot be published without one. The \texttt{reset\_at} column is indexed
because the take screen queries it on every load to invalidate stale cached
sessions when a teacher destructively edits an exam.

The single most interesting table in the schema is \texttt{Problem}, which is
both \emph{polymorphic} and \emph{dual-purpose}. It is polymorphic because
one table holds all three question types --- coding, MCQ, and written ---
discriminated by a \texttt{type} column (\texttt{size:20;not
null;default:coding}). Rather than three normalized tables joined at query
time, every variant-specific column lives on one wide row and is simply left
null for the types that do not use it:

\begin{itemize}[nosep]
  \item \textbf{Coding} problems use \texttt{starter\_code},
        \texttt{time\_limit\_ms} (default 2000), and
        \texttt{memory\_limit\_kb} (default 262144, i.e.\ 256\,MB), plus
        child \texttt{TestCase} rows.
  \item \textbf{MCQ} problems use \texttt{options} and
        \texttt{correct\_option\_ids} (both JSONB), \texttt{multiple\_correct},
        and \texttt{explanation}.
  \item \textbf{Written} problems use \texttt{max\_word\_count} (default 500),
        \texttt{rubric}, and \texttt{require\_manual\_grading}.
\end{itemize}

The trade-off is explicit and, at this scale, defensible: a wider table with
more nullable columns is exchanged for a much simpler grading dispatch. The
backend performs a single \texttt{switch Problem.Type} rather than joining
three type-specific tables, and the exam's ``problem list'' is one
homogeneous query~\cite{fowler2002patterns}. It is dual-purpose because
\texttt{ExamID} is a nullable \texttt{*uint} and \texttt{IsBank} is a
boolean: an ordinary problem has \texttt{ExamID} set and
\texttt{is\_bank=false}, while a reusable bank template has
\texttt{ExamID = NULL} and \texttt{is\_bank=true}, unbound to any exam and
optionally filed under a \texttt{Folder}.

Each coding \texttt{Problem} owns many \texttt{TestCase} rows. The crucial
column there is \texttt{is\_sample}: sample cases are the ones surfaced to
students through the live ``Run Samples'' feature, while hidden cases are used
only by the real auto-grader and are never serialized to the client. This
split is the data side of a defense-in-depth boundary --- hidden expected
outputs never leave the server.

\subsection{JSONB Columns}\label{sec:jsonb}

APEX uses PostgreSQL's binary-JSON type, JSONB, in four places:
\texttt{problems.options} and \texttt{problems.correct\_option\_ids} (MCQ
option sets), \texttt{problems.tags} (default \texttt{'[]'}),
\texttt{announcements.attachments}, and \texttt{exam\_attempts.draft\_answers}
(discussed in \cref{sec:attempt-aggregate}). The common thread is that each
holds a small, variable-length, client-shaped blob that the application reads
and writes \emph{whole} and never needs to query \emph{inside} at the SQL
level. Kleppmann observes that document-style storage is a good fit exactly
when data has a natural tree shape that is loaded in its entirety and rarely
joined against~\cite{kleppmann2017ddia}; that is precisely the case here. The
\texttt{draft\_answers} payload in particular must round-trip the client's
exact structure, so storing a separate normalized table would add joins for
no query benefit. The cost --- that these fields are opaque to SQL
predicates --- is one APEX never pays, because it never filters on their
interior.

\section{The Attempt Aggregate}\label{sec:attempt-aggregate}

The taking half of the assessment cluster is the correctness core of the
whole system. Its shape is a two-level ownership chain:
\texttt{ExamAttempt} $\rightarrow$ (1:N) $\rightarrow$ \texttt{Submission}
$\rightarrow$ (1:N) $\rightarrow$ \texttt{TestResult}. One sitting
(\texttt{ExamAttempt}) collects one \texttt{Submission} per answered problem;
each coding submission collects one \texttt{TestResult} per test case Judge0
ran against it.

\begin{lstlisting}[language=Go,caption={The ExamAttempt root and its unique index (internal/models/exam\_attempt.go).},label={lst:attempt}]
type ExamAttempt struct {
    ID          uint       `gorm:"primaryKey" json:"id"`
    UserID      uint       `gorm:"not null;index;uniqueIndex:idx_user_exam_attempt" json:"user_id"`
    ExamID      uint       `gorm:"not null;index;uniqueIndex:idx_user_exam_attempt" json:"exam_id"`
    StartedAt   time.Time  `gorm:"autoCreateTime" json:"started_at"`
    SubmittedAt *time.Time `json:"submitted_at,omitempty"`
    Score       float64    `gorm:"not null;default:0" json:"score"`
    Status      string     `gorm:"size:30;not null;default:in_progress" json:"status"`
    GradedNotified bool    `gorm:"not null;default:false" json:"graded_notified"`
    DraftAnswers   *string    `gorm:"type:jsonb" json:"draft_answers,omitempty"`
    DraftSavedAt   *time.Time `json:"draft_saved_at,omitempty"`
    Submissions []Submission `gorm:"foreignKey:ExamAttemptID" json:"submissions,omitempty"`
}
\end{lstlisting}

The load-bearing line is the third composite unique index,
\texttt{idx\_user\_exam\_attempt}, declared on both \texttt{UserID} and
\texttt{ExamID}. It enforces \textbf{one attempt row per (student, exam)
pair}, and that single constraint is what makes two otherwise-hard features
tractable.

\paragraph{Resume-after-crash.} Because the pair is unique, the database
physically cannot hold two attempt rows for the same student and exam. The
\texttt{StartAttempt} handler is therefore idempotent: it looks up the
attempt by \texttt{(user\_id, exam\_id)}, returns the existing row if one is
found, and creates a row only when none exists. A student who refreshes the
page, whose laptop dies, or who reopens the exam on a second device resolves
every time to the \emph{same} attempt. In-progress answers survive because
the autosave endpoint writes the raw answer body into the
\texttt{draft\_answers} JSONB column of that same attempt, stamping
\texttt{draft\_saved\_at}; on reconnect the take screen rehydrates the draft
and recomputes remaining time from \texttt{started\_at}. There is exactly one
place the draft can live, which is what ties autosave to resume.

\paragraph{Safe re-scoring.} The attempt's \texttt{score} is stored on the
row but \emph{derived by recomputation}, never incrementally mutated. The
finalisation routine iterates the exam's \emph{full} problem list ---
so a skipped question stays in the denominator at zero and a student who
answered only the easy questions cannot appear at 100\,\% --- sums earned
points over total points, and writes back a percentage. Written answers in
\texttt{pending\_review} are excluded from both numerator and denominator
until a teacher grades them, so an in-flight manual grade does not depress
the visible total. Because there is exactly one attempt, and the score is a
pure re-sum over a fixed set of children, the routine is safe to call any
number of times: after each test result, after a manual grade, after a retry,
the same inputs always yield the same number. If duplicate attempts were
possible, this re-sum could group the wrong submissions or double-count; the
unique index removes that entire failure class.

\paragraph{Idempotent state transitions.} The attempt's \texttt{status}
moves \texttt{in\_progress} $\rightarrow$ \texttt{submitted}. The
\texttt{graded\_notified} boolean is an idempotency guard so the ``Exam
Graded'' notification fires exactly once even though finalisation may run
many times; the finaliser sets it under the ``all graded, nothing in
\texttt{pending\_review}'' condition. This is the same philosophy as the
reminder-dedup timestamp columns on \texttt{Exam} --- idempotency expressed
as a database column rather than in-memory state, so a process restart can
never double-fire.

Below the attempt, \texttt{Submission} carries the student's code, selected
options (JSONB), or text answer, along with its own grading verdict
(\texttt{status}, \texttt{passed\_count}, \texttt{total\_count},
\texttt{score}, and the maximum execution time and memory observed across its
tests). Its \texttt{exam\_attempt\_id} is a nullable \texttt{*uint}: standalone
playground and sample runs create no attempt, so the pointer tolerates that
context while every real exam submission is grouped under its attempt.
\texttt{TestResult}, the leaf, records one row per test case: whether it
passed, the actual output, timing, memory, and the raw Judge0 status.

\Cref{lst:examattempts-ddl} shows the \texttt{exam\_attempts} table as
\texttt{pg\_dump} emits it, with the unique constraint spelled out in SQL.

\begin{lstlisting}[language=SQL,caption={The exam\_attempts table, from the schema dump (Appendix~B).},label={lst:examattempts-ddl}]
CREATE TABLE exam_attempts (
    id              BIGSERIAL PRIMARY KEY,
    user_id         BIGINT NOT NULL,
    exam_id         BIGINT NOT NULL,
    started_at      TIMESTAMPTZ DEFAULT now(),
    submitted_at    TIMESTAMPTZ,
    score           DOUBLE PRECISION NOT NULL DEFAULT 0,
    status          VARCHAR(30) NOT NULL DEFAULT 'in_progress',
    graded_notified BOOLEAN NOT NULL DEFAULT false,
    draft_answers   JSONB,
    draft_saved_at  TIMESTAMPTZ,
    UNIQUE (user_id, exam_id)
);
\end{lstlisting}

\section{Indexes and Soft Constraints}\label{sec:indexes}

APEX leans on the database for the invariants it cannot afford to get wrong.
Three composite unique indexes, summarized in \cref{tab:indexes}, each enforce
a ``no duplicate pair'' rule at the storage layer, so the guarantee holds even
under a race between two concurrent requests --- the second insert simply
fails with a unique-violation, and at most one row can ever exist.

\begin{table}[htbp]
\centering
\caption{The three composite unique indexes and their invariants.}
\label{tab:indexes}
\begin{tabularx}{\linewidth}{@{}L{3.6cm} L{3.4cm} X@{}}
\toprule
\textbf{Index} & \textbf{Columns} & \textbf{Invariant enforced} \\
\midrule
\texttt{idx\_class\_user}        & \texttt{class\_members\-(class\_id, user\_id)}   & A student joins a given class at most once. \\
\texttt{idx\_user\_exam\_attempt} & \texttt{exam\_attempts\-(user\_id, exam\_id)}    & Exactly one attempt row per (student, exam) --- the linchpin of resume and re-scoring. \\
\texttt{idx\_exam\_class}        & \texttt{exam\_classes\-(exam\_id, class\_id)}    & An exam is assigned to a given class at most once. \\
\bottomrule
\end{tabularx}
\end{table}

The reason \texttt{idx\_user\_exam\_attempt} matters more than the other two
is that it is not merely a hygiene constraint against duplicate rows; it is
the foundation of a behavioural guarantee. Enforcing ``one attempt per
student per exam'' in the database rather than in application code means the
resume and re-scoring logic never has to reason about ``which attempt is the
real one.'' There is provably only one. Had the team allowed multiple
attempts and picked ``the latest'' in code, that correctness rule would have
had to be threaded through every query touching attempts, and a refresh storm
could have created two rows the finaliser would then aggregate incorrectly.
Putting the rule in the schema collapses a whole category of race conditions
into a single \texttt{UNIQUE} clause. The honest cost is that APEX supports
exactly one attempt per exam with no native retake; adding retakes later
would mean evolving the constraint to \texttt{(user\_id, exam\_id,
attempt\_number)} --- a deliberate schema change, not an accident.

A fourth index deserves mention because it shows the limit of what a GORM
struct tag can express. The \texttt{google\_id} column is declared as a plain
\texttt{index} in the struct, but true uniqueness is enforced by a
\emph{partial} unique index built in hand-written SQL:

\begin{lstlisting}[language=SQL,caption={The partial unique index on google\_id (internal/database/migrations.go).},label={lst:google-idx}]
CREATE UNIQUE INDEX IF NOT EXISTS idx_users_google_id
    ON users(google_id)
    WHERE google_id IS NOT NULL AND google_id <> '';
\end{lstlisting}

Most users are password-only and have an empty \texttt{google\_id}; a plain
unique index would treat all those empty strings as colliding duplicates. The
partial index applies uniqueness \emph{only} to present, non-empty Google
subjects, so password users never collide while Google-linked accounts stay
uniquely bound. GORM's tag vocabulary cannot express a conditional index,
which is one concrete reason the project keeps a hand-written SQL layer at
all.

Beyond the unique indexes, every hot-path foreign key
(\texttt{user\_id}, \texttt{exam\_id}, \texttt{exam\_attempt\_id},
\texttt{submission\_id}, \texttt{problem\_id}, \texttt{reset\_at}) carries a
plain \texttt{index} tag so lookups stay logarithmic. What APEX largely does
\emph{not} do is declare database-level referential foreign-key
\emph{constraints} between most tables. Relationships are expressed as GORM
associations and enforced by application logic --- ``soft'' constraints ---
rather than \texttt{REFERENCES} clauses with \texttt{ON DELETE CASCADE}. The
sole hard foreign key added deliberately is \texttt{fk\_problems\_folder}
with \texttt{ON DELETE SET NULL}, because folder deletion has a specific
non-cascading semantics (unfile, do not destroy). The trade-off is honest:
soft constraints keep AutoMigrate simple and account deletion tractable (the
cascade is coded explicitly in a transaction rather than left to the
database), at the price of relying on the application to keep references
consistent. Account deletion is precisely where this shows --- it is a
single hand-written transactional cascade that clears profile,
notifications, reset tokens, and role-scoped children in order before
deleting the user.

\section{The Migration Pipeline}\label{sec:migrations}

A schema is never finished; columns get added, defaults change, data gets
backfilled. APEX runs two parallel migration tracks, and which one mutates the
schema is decided by a single environment flag, \texttt{SKIP\_AUTOMIGRATE},
read in \texttt{internal/database/database.go}. \Cref{fig:migration} shows the
decision and the three-stage development path.

\begin{figure}[htbp]
\centering
\resizebox{\textwidth}{!}{%
\begin{tikzpicture}[
  node distance=6mm,
  font=\sffamily\scriptsize,
  mstep/.style={apexbox, text width=30mm, minimum height=11mm},
  legacy/.style={apexstore, text width=34mm, minimum height=11mm, align=left, font=\sffamily\scriptsize},
  auth/.style={apexext, text width=32mm, minimum height=11mm},
  mdec/.style={draw=apexblue, fill=apexlight, rounded corners=2pt, align=center,
    font=\sffamily\scriptsize\bfseries, minimum width=20mm, minimum height=9mm},
]

% --- Boot entry ---
\node[mstep] (boot) {\textbf{Process boot}\\\texttt{main.go}\\\texttt{database.Init}};

% --- Idempotent legacy SQL pass ---
\node[legacy, right=10mm of boot] (legacy) {%
  \textbf{Idempotent legacy SQL}\\\texttt{RunMigrations(DB)}\\[2pt]
  \footnotesize
  1. \texttt{ADD COLUMN} nullable\\
  2. backfill correct value\\
  3. \texttt{SET NOT NULL}\\
  4. \texttt{SET DEFAULT}\\
  5. repair rows\\[1pt]
  \emph{each block gated by} \texttt{tableExists}};

% --- Decision on SKIP_AUTOMIGRATE ---
\node[mdec, right=11mm of legacy] (skip) {\texttt{SKIP\_}\\\texttt{AUTOMIGRATE?}};

% --- Dev branch: GORM AutoMigrate ---
\node[mstep, above right=6mm and 12mm of skip, text width=34mm] (auto) {%
  \textbf{GORM \texttt{AutoMigrate}}\\
  new tables only, all 15 models\\
  then \texttt{RunPostMigrations}\\
  \emph{(dev / unset $=$ false)}};

% --- Prod branch: golang-migrate ---
\node[mstep, below right=6mm and 12mm of skip, text width=34mm,
  draw=apexgreen!70!black, fill=apexgreen!8] (migrate) {%
  \textbf{Versioned \texttt{golang-migrate}}\\
  \texttt{cmd/migrate up}\\
  \texttt{migrations/sql/NNNN\_*}\\
  \emph{SOLE schema authority (prod, }\texttt{true}\emph{)}};

% --- Arrows ---
\draw[apexarrowblue] (boot) -- (legacy);
\draw[apexarrowblue] (legacy) -- (skip);
\draw[apexarrow] (skip.north) to[out=90,in=180] node[apexlabel, pos=0.55] {false} (auto.west);
\draw[apexarrow] (skip.south) to[out=-90,in=180] node[apexlabel, pos=0.55] {true} (migrate.west);

% --- is_draft incident callout ---
\node[apexext, below=13mm of legacy, text width=0.92\textwidth, align=left,
  draw=apexamber, line width=0.8pt, font=\sffamily\scriptsize] (incident) {%
  \textbf{\color{apexamber}The \texttt{is\_draft} incident (war story).}\;
  A \texttt{gorm:"default:true"} tag on \texttt{Exam.IsDraft} made \texttt{AutoMigrate}
  apply that \emph{DDL default to every existing row}, flipping all exams to
  \texttt{is\_draft = true}. Student lists filter \texttt{WHERE is\_draft = false},
  so every exam \emph{silently vanished} — no error logged.\\[2pt]
  \textbf{Fix:} (a) struct default flipped to safe zero-value \texttt{false} (``published'');\;
  (b) an explicit \emph{idempotent} SQL block owns the change
  (\texttt{ADD COLUMN IF NOT EXISTS} $\rightarrow$ backfill $\rightarrow$
  \texttt{SET NOT NULL} $\rightarrow$ \texttt{SET DEFAULT false} $\rightarrow$ repair);\;
  (c) production runs with \texttt{AutoMigrate} off (\texttt{SKIP\_AUTOMIGRATE=true}),
  leaving \texttt{cmd/migrate} the sole schema authority.};

% link the incident to the AutoMigrate node it originated from
\draw[apexarrow, dashed, draw=apexamber]
  (auto.south) to[out=-90,in=90] node[apexlabel, pos=0.5] {silent corruption} (incident.north east);

\end{tikzpicture}%
}
\caption{Schema-migration pipeline and the is\_draft incident.}
\label{fig:migration}
\end{figure}


\subsection{The Development Path: Three Stages, One Direction}

When \texttt{SKIP\_AUTOMIGRATE} is unset or false (the developer laptop),
\texttt{Connect()} opens the pool --- 25 max-open connections, 10 idle, a
five-minute idle lifetime and a thirty-minute maximum lifetime --- and then
runs three passes in a fixed order.

\textbf{Stage 1 --- \texttt{RunMigrations(DB)}}, hand-written idempotent SQL,
runs \emph{before} AutoMigrate. This is where all the work that AutoMigrate
cannot do safely lives: data backfills, nullability changes, default changes,
the partial Google index, dropping legacy tables, and the historical score
re-aggregation. Every block is gated on the target table already existing, via
a helper:

\begin{lstlisting}[language=Go,caption={The table-existence gate (internal/database/migrations.go).},label={lst:tableexists}]
func tableExists(db *gorm.DB, name string) bool {
    var n int64
    db.Raw("SELECT 1 FROM information_schema.tables "+
        "WHERE table_schema = current_schema() AND table_name = ?", name).Scan(&n)
    return n > 0
}
\end{lstlisting}

On a first boot against an empty database, every gate is false, so every
block is skipped and Stage~2 creates the tables for real. On the second boot
the gates are true and the SQL replays harmlessly, because every statement is
written with \texttt{IF NOT EXISTS} / \texttt{IF EXISTS} guards. Idempotency
here is deliberate: an operator restarting a misbehaving process at three in
the morning should never have to reason about migration ordering. As the
team's own lesson puts it, idempotency is cheaper than coordination.

\textbf{Stage 2 --- \texttt{AutoMigrate(...)}}, GORM's automatic schema sync,
creates only brand-new tables, columns, and indexes for the fifteen
registered models. It does the easy, safe part; the dangerous part already
happened in Stage~1.

\textbf{Stage 3 --- \texttt{RunPostMigrations(DB)}} runs \emph{after}
AutoMigrate because it needs a table that AutoMigrate just created. It sweeps
orphaned \texttt{folder\_id} references to null, then adds the
\texttt{fk\_problems\_folder} foreign key referencing \texttt{folders}. The
ordering is forced by dependency: you cannot add a foreign key to a table that
does not yet exist, and you cannot enforce the constraint while dangling rows
violate it, so the orphan sweep precedes the \texttt{ADD CONSTRAINT}. The
whole pipeline moves in one direction --- forward only.

\subsection{The Production Path: A Versioned Binary Owns the Schema}

When \texttt{SKIP\_AUTOMIGRATE=true}, \texttt{Connect()} takes an entirely
different branch: it calls \texttt{RunSQLMigrations(cfg)} and returns,
\emph{never} touching \texttt{RunMigrations}, \texttt{AutoMigrate}, or
\texttt{RunPostMigrations}. \texttt{RunSQLMigrations} drives the
\texttt{golang-migrate} library~\cite{golangmigrate} over the numbered,
paired \texttt{NNNN\_name.up.sql} / \texttt{.down.sql} files in
\texttt{internal/database/migrations/sql/} (directory overridable via
\texttt{MIGRATIONS\_DIR}). It handles \texttt{ErrNoChange} and
\texttt{ErrNilVersion} gracefully and logs the current version and dirty flag.
The same logic is exposed as a standalone \texttt{cmd/migrate} CLI supporting
\texttt{up}, \texttt{down}, \texttt{version}, and \texttt{force~<n>}, which the
production Docker stack runs as a one-shot \texttt{migrate} service before the
backend starts.

Today the versioned track contains exactly one pair, the baseline
\texttt{0001\_init}, which is an intentional no-op (\texttt{SELECT 1;}) whose
sole purpose is to mark the AutoMigrated schema as version~1 so that future
\texttt{0002+} files can layer explicit, reviewed changes on top. The reason
production skips AutoMigrate entirely is a matter of principle earned through
experience: AutoMigrate's schema \emph{inference} is a convenience on a
laptop, but a production database wants deterministic, version-controlled,
reversible DDL that a human has reviewed --- not an ORM silently reshaping a
populated table~\cite{twelvefactor}. That silent reshaping is exactly what
caused the incident described next.

\section{Case Study: The \texttt{is\_draft} Incident}\label{sec:incident}

The clearest justification for the entire migration discipline is a single
bug the team shipped, caught, and turned into a written rule. It is worth
telling in full because it is a textbook example of how an ORM can do
\emph{more} than the programmer intended.

\paragraph{The mechanism.} Someone added the tag \texttt{gorm:"default:true"}
to the \texttt{IsDraft} field on the \texttt{Exam} struct, intending it to
mean ``new exams start as drafts.'' On the next boot, GORM AutoMigrate
translated that tag into a DDL default and applied it to the
\emph{already-populated} \texttt{exams} table, flipping every existing exam
row to \texttt{is\_draft = true}. The trouble is that student-facing list
queries filter \texttt{WHERE is\_draft = false} --- students should see only
published exams. With every row now marked draft, that filter matched
nothing. \textbf{Every exam silently vanished from every student's view.}
Nothing crashed, nothing was logged; the data simply became invisible. This
is the definition of silent corruption, and it is the failure mode that makes
AutoMigrate dangerous on populated tables: GORM adds the column and applies
the default, but it has no way to know that the ``correct'' value for
existing rows is the opposite of the new default~\cite{gorm}.

\paragraph{The multi-part fix.} The repair had three independent parts, each
addressing a different layer.

First, the struct default was flipped to the safe zero-value:

\begin{lstlisting}[language=Go,caption={The corrected field: the safe zero-value is ``published.''},label={lst:isdraft-fix}]
IsDraft bool `gorm:"default:false" json:"is_draft"`
\end{lstlisting}

Now the default that AutoMigrate might apply is \texttt{false} ---
``published'' --- so even a stray application is harmless. Keeping the tag is
legitimate: it governs the INSERT-time default when the Go field is
zero-valued, which is a safe, application-level use.

Second, and more importantly, the real correctness was moved into explicit,
idempotent SQL in \texttt{RunMigrations}, following the canonical add-nullable
$\rightarrow$ backfill $\rightarrow$ constrain pattern, plus a repair line for
rows the bad migration had already corrupted:

\begin{lstlisting}[language=SQL,caption={The idempotent fix that owns the is\_draft column (internal/database/migrations.go).},label={lst:isdraft-sql}]
-- Step 1: add the column WITHOUT forcing a default onto existing rows
ALTER TABLE exams ADD COLUMN IF NOT EXISTS is_draft BOOLEAN;
-- Step 2: backfill existing rows with the CORRECT value
UPDATE exams SET is_draft = false WHERE is_draft IS NULL;
-- Step 3: only now is it safe to constrain
ALTER TABLE exams ALTER COLUMN is_draft SET NOT NULL;
ALTER TABLE exams ALTER COLUMN is_draft SET DEFAULT false;
-- Repair: rows the bad migration flipped, but which clearly WERE published
UPDATE exams SET is_draft = false
    WHERE is_draft = true AND start_time IS NOT NULL;
\end{lstlisting}

The repair line uses \texttt{start\_time IS NOT NULL} as a reliable signal
that an exam was genuinely published, because a draft cannot be published
without a start time. Adding the column nullable first is the whole point: no
default is forced onto existing rows, so the backfill in Step~2 can set the
\emph{correct} value before the \texttt{NOT NULL} and \texttt{DEFAULT}
constraints in Step~3.

Third, production was configured to run with \texttt{SKIP\_AUTOMIGRATE=true},
making the versioned \texttt{migrate} binary the sole authority over the
production schema. GORM can never again reshape a populated production table,
because in production GORM never runs DDL at all.

\paragraph{The resulting rule.} The incident produced \texttt{ORM\_RULES.md},
a short file of hard rules at the repository root. Its first and central rule
is: \emph{never use \texttt{gorm:"default:X"} on existing tables}. GORM
AutoMigrate only \emph{adds} columns; it does not manage defaults safely on
populated tables, so existing rows may receive the default unexpectedly. The
\texttt{gorm:"default"} tag is reserved for application-level INSERT defaults
only, never for DDL backfill. The supporting rules follow directly: audit
every query when adding a filter column (the \texttt{WHERE is\_draft = false}
clause was the amplifier that turned a bad default into an outage); define the
zero-value meaning of every new boolean explicitly (what does \texttt{false}
mean, what should existing rows be); write existing-table changes as explicit
idempotent SQL run \emph{before} AutoMigrate; and verify row counts after any
migration before touching application code. The fix is, in the end, exactly
the shape of those rules --- and the three-stage pipeline of
\cref{sec:migrations} is their operational form.

\section{Design Trade-offs and Honest Gaps}

The data layer's choices were made in favour of smallness and provable
correctness over normalization purity and speculative flexibility, and it is
worth naming the costs honestly.

\begin{itemize}[nosep]
  \item \textbf{The polymorphic \texttt{problems} table is wide}, with columns
        that are always null for two of the three question types. This is the
        price of a one-line grading dispatch; a table-per-type design would be
        cleaner on paper but would trade a homogeneous problem-list query for
        three-way joins.
  \item \textbf{One attempt per exam, no retakes.} The unique index that makes
        resume and re-scoring tractable also fixes the model at a single
        attempt. Retakes are out of scope for this version, and the constraint
        makes adding them a deliberate schema evolution rather than an
        accident.
  \item \textbf{Two migration tracks coexist.} The legacy Go-SQL
        \texttt{RunMigrations} still owns the baseline locally, while the
        versioned \texttt{golang-migrate} track is the production future
        bridged by the no-op \texttt{0001\_init}. This is transitional and
        documented, but it does mean two places to look when adding a change.
  \item \textbf{Soft constraints over hard foreign keys.} Most relationships
        are enforced by application logic, not database \texttt{REFERENCES}.
        This keeps AutoMigrate and account-deletion cascades simple at the
        cost of relying on the application to keep references consistent.
  \item \textbf{Single-instance assumptions.} The idempotent dedup columns
        make restarts safe, but they are not a substitute for a durable work
        queue; the fire-and-forget grading path and the reminder scheduler
        both assume a single backend process.
\end{itemize}

The through-line of this chapter is that APEX treats the database not as a
passive bucket but as the place where the hardest invariants are enforced.
The unique \texttt{(user\_id, exam\_id)} index is the single fact the entire
grading-correctness story rests on; the migration pipeline exists so that the
schema can evolve without ever again silently violating an invariant like it
did the day every exam vanished. Small, auditable, and disciplined about
change --- that is the whole design.
